\documentclass[10pt]{beamer}

% ============================================================
% CLASSICAL ECONOMICS BEAMER TEMPLATE
% Clean academic style for seminars, field courses, job talks,
% conference presentations, and paper presentations.
%
% QUICK CHANGES
% 1. For widescreen, replace the first line with:
%    \documentclass[10pt,aspectratio=169]{beamer}
% 2. Change title/author/institute/date below.
% 3. Delete any slide you do not need.
% ============================================================

% -------------------- Theme --------------------
\usetheme{Boadilla}
\usefonttheme{serif}
\setbeamertemplate{navigation symbols}{}
\setbeamertemplate{caption}[numbered]

% Classical, restrained colors
\definecolor{econblue}{RGB}{28,63,105}
\definecolor{econgray}{RGB}{90,90,90}
\setbeamercolor{structure}{fg=econblue}
\setbeamercolor{title}{fg=econblue}
\setbeamercolor{frametitle}{fg=econblue}
\setbeamercolor{footline}{fg=econgray,bg=white}

% Footer: author | short title | page number
\setbeamertemplate{footline}{%
  \leavevmode%
  \hbox{%
  \begin{beamercolorbox}[wd=.32\paperwidth,ht=2.5ex,dp=1.1ex,leftskip=.8em]{author in head/foot}%
    \usebeamerfont{author in head/foot}\insertshortauthor
  \end{beamercolorbox}%
  \begin{beamercolorbox}[wd=.50\paperwidth,ht=2.5ex,dp=1.1ex,center]{title in head/foot}%
    \usebeamerfont{title in head/foot}\insertshorttitle
  \end{beamercolorbox}%
  \begin{beamercolorbox}[wd=.18\paperwidth,ht=2.5ex,dp=1.1ex,rightskip=.8em plus1fil]{date in head/foot}%
    \insertframenumber
  \end{beamercolorbox}}%
  \vskip0pt%
}

% -------------------- Packages --------------------
\usepackage{amsmath,amssymb,mathtools,bm}
\usepackage{booktabs}
\usepackage{threeparttable}
\usepackage{graphicx}
\usepackage{array}
\usepackage{siunitx}
\usepackage{tikz}
\usepackage{appendixnumberbeamer}

% -------------------- Small helper commands --------------------
\newcommand{\E}{\mathbb{E}}
\newcommand{\Var}{\operatorname{Var}}
\newcommand{\Cov}{\operatorname{Cov}}

% -------------------- Title information --------------------
\title[Short Paper Title]{A Classical Economics Presentation Title}
\subtitle{Optional Subtitle or Paper Version}
\author[Your Name]{Your Name}
\institute[University]{Department of Economics \\ Your University}
\date{Seminar / Course / Conference \\ September 2026}

% ============================================================
\begin{document}
% ============================================================

% -------------------- Title slide --------------------
\begin{frame}
  \titlepage
\end{frame}

% Optional: remove this slide for short talks.
\begin{frame}{Roadmap}
  \tableofcontents
\end{frame}

% ============================================================
\section{Motivation}
% ============================================================

\begin{frame}{Motivation}
\begin{itemize}
  \item Start with the economic question, not the method.
  \item Give one or two facts that make the question important.
  \item Explain the gap in the existing literature.
  \item End with what your paper contributes.
\end{itemize}

\vspace{0.6em}
\begin{block}{Main question}
How does $X$ affect $Y$, and through which economic mechanism?
\end{block}
\end{frame}

\begin{frame}{What This Paper Does}
\begin{enumerate}
  \item Introduces a new setting / dataset / source of variation.
  \item Estimates the main relationship of interest.
  \item Studies mechanisms or heterogeneous effects.
  \item Connects the results to a broader economic question.
\end{enumerate}

\vspace{0.8em}
\textbf{Main takeaway:} State the result in one sentence with an economically meaningful magnitude.
\end{frame}

% ============================================================
\section{Background and Framework}
% ============================================================

\begin{frame}{Institutional Background}
\begin{columns}[T]
\column{0.56\textwidth}
\begin{itemize}
  \item What is the institution, policy, market, or environment?
  \item What should the audience know before seeing the empirical design?
  \item Which dates, rules, or thresholds matter?
\end{itemize}

\column{0.40\textwidth}
\begin{block}{Keep this slide selective}
Only include institutional details that are needed to understand the economic mechanism or research design.
\end{block}
\end{columns}
\end{frame}

\begin{frame}{A Simple Economic Framework}
A generic relationship can be written as
\[
  y_i = f(x_i,z_i;\theta) + \varepsilon_i,
  \qquad i=1,\ldots,n.
\]

\begin{itemize}
  \item $y_i$: outcome of interest.
  \item $x_i$: key explanatory variable.
  \item $z_i$: other observed characteristics.
  \item $\theta$: parameter(s) with economic interpretation.
\end{itemize}

\vspace{0.4em}
\alert{Presentation tip:} do not put every equation from the paper on slides. Show only equations needed for the argument.
\end{frame}

% Example of aligned equations
\begin{frame}{Useful Equation Layout}
\begin{align}
  U(c_1,c_2) &= u(c_1) + \beta u(c_2), \\
  c_1 + \frac{c_2}{1+r} &= w.
\end{align}

\begin{itemize}
  \item Put alignment marker \texttt{\&} before the equality sign.
  \item Use \texttt{\\} to start the next aligned line.
  \item Use \texttt{\qquad} for a clean horizontal separation inside math.
\end{itemize}
\end{frame}

% ============================================================
\section{Data and Research Design}
% ============================================================

\begin{frame}{Data}
\begin{itemize}
  \item \textbf{Source:} Administrative / survey / transaction / historical data.
  \item \textbf{Unit of observation:} Individual, firm, county, country, etc.
  \item \textbf{Period:} 2000--2025.
  \item \textbf{Main outcome:} Define it in plain English first.
  \item \textbf{Sample:} Explain key restrictions.
\end{itemize}

\vspace{0.6em}
\begin{block}{Good data slide rule}
The audience should know what one row of your dataset represents before you show a regression table.
\end{block}
\end{frame}

\begin{frame}{Summary Statistics}
\centering
\small
\begin{threeparttable}
\begin{tabular}{lS[table-format=4.2]S[table-format=4.2]S[table-format=5.0]}
\toprule
Variable & {Mean} & {Std. Dev.} & {$N$} \\
\midrule
Outcome $Y$      & 12.40 & 4.10 & 12500 \\
Treatment $X$    & 0.36  & 0.48 & 12500 \\
Income           & 58.20 & 21.70 & 12500 \\
Age              & 41.30 & 12.60 & 12500 \\
\bottomrule
\end{tabular}
\begin{tablenotes}\footnotesize
\item \textit{Notes:} Replace this with a concise description of the sample and units.
\end{tablenotes}
\end{threeparttable}
\end{frame}

\begin{frame}{Research Design}
\begin{block}{Core idea}
Explain the source of identifying variation in words before showing the estimating equation.
\end{block}

\vspace{0.6em}
A generic specification is
\[
  Y_{it} = \alpha + \beta X_{it} + \gamma' Z_{it} + \mu_i + \tau_t + \varepsilon_{it}.
\]

\begin{itemize}
  \item $\beta$ is the parameter of interest.
  \item $\mu_i$ denotes unit fixed effects.
  \item $\tau_t$ denotes time fixed effects.
\end{itemize}
\end{frame}

% ============================================================
\section{Results}
% ============================================================

\begin{frame}{Main Result}
\centering
\small
\begin{threeparttable}
\begin{tabular}{lccc}
\toprule
 & (1) & (2) & (3) \\
\midrule
$X$ & 0.118*** & 0.104*** & 0.072** \\
    & (0.031)  & (0.034)  & (0.029) \\
\addlinespace
Controls      & No  & Yes & Yes \\
Unit FE       & No  & No  & Yes \\
Time FE       & No  & Yes & Yes \\
Observations  & 12,500 & 12,500 & 12,500 \\
$R^2$         & 0.11 & 0.24 & 0.39 \\
\bottomrule
\end{tabular}
\begin{tablenotes}\scriptsize
\item \textit{Notes:} Standard errors in parentheses. *, **, and *** denote significance at the 10\%, 5\%, and 1\% levels.
\end{tablenotes}
\end{threeparttable}

\vspace{0.7em}
\textbf{Interpretation:} translate the preferred coefficient into an economically meaningful magnitude.
\end{frame}

\begin{frame}{Main Figure}
\centering

% This is a self-contained placeholder figure made with TikZ.
% Replace it with:
% \includegraphics[width=0.78\linewidth]{figures/main_figure.pdf}
\begin{tikzpicture}[x=0.75cm,y=0.65cm]
  \draw[->] (0,0) -- (8.3,0) node[right] {Period};
  \draw[->] (0,-2.3) -- (0,2.5) node[above] {Estimate};
  \draw[dashed,gray] (0,0) -- (8,0);
  \draw[thick,econblue] plot[smooth] coordinates {
    (0,-0.6) (1,-0.3) (2,-0.1) (3,0.0) (4,0.5) (5,0.9) (6,1.25) (7,1.45)
  };
  \foreach \x/\y in {0/-0.6,1/-0.3,2/-0.1,3/0.0,4/0.5,5/0.9,6/1.25,7/1.45}
    \fill[econblue] (\x,\y) circle (1.7pt);
  \draw[densely dashed] (3,-2.1) -- (3,2.2);
  \node[below] at (3,0) {0};
\end{tikzpicture}

\vspace{0.5em}
\begin{minipage}{0.85\textwidth}
\footnotesize
\textit{Notes:} Keep figure notes short on slides. Explain the key pattern orally and use a clear title that states the result.
\end{minipage}
\end{frame}

\begin{frame}{How to Present a Result}
\begin{enumerate}
  \item State the sign and direction.
  \item Give the magnitude.
  \item Translate the magnitude into economically meaningful units.
  \item Say what comparison or variation identifies it.
  \item Then discuss statistical precision.
\end{enumerate}

\vspace{0.7em}
\begin{alertblock}{Avoid}
Reading every coefficient and standard error from the table.
\end{alertblock}
\end{frame}

% ============================================================
\section{Mechanisms and Robustness}
% ============================================================

\begin{frame}{Mechanism}
\begin{columns}[T]
\column{0.47\textwidth}
\begin{block}{Mechanism 1}
Economic channel and corresponding empirical prediction.
\end{block}

\column{0.47\textwidth}
\begin{block}{Mechanism 2}
Alternative channel and how the data distinguish it.
\end{block}
\end{columns}

\vspace{1em}
\textbf{Goal:} connect the mechanism tests to the conceptual framework, rather than presenting disconnected regressions.
\end{frame}

\begin{frame}{Robustness}
\begin{itemize}
  \item Alternative sample definitions.
  \item Alternative outcome or treatment definitions.
  \item Alternative controls / fixed effects.
  \item Placebo or falsification tests.
  \item Alternative inference / clustering choices.
\end{itemize}

\vspace{0.7em}
\begin{block}{Presentation tip}
Do not show ten robustness tables in the main talk. Summarize the pattern and move detailed tables to the appendix.
\end{block}
\end{frame}

% ============================================================
\section{Conclusion}
% ============================================================

\begin{frame}{Conclusion}
\begin{enumerate}
  \item \textbf{Question:} one-line reminder of the research question.
  \item \textbf{Finding:} main empirical or theoretical result.
  \item \textbf{Mechanism:} what explains the result.
  \item \textbf{Implication:} why the finding matters for economics.
\end{enumerate}

\vspace{1em}
\centering
\Large Thank you
\end{frame}

% ============================================================
% APPENDIX / BACKUP SLIDES
% Slides after \appendix are usually backup slides.
% appendixnumberbeamer prevents them from changing the main slide count.
% ============================================================
\appendix
\section{Appendix}

\begin{frame}{Appendix: Additional Result}
\begin{itemize}
  \item Put detailed robustness checks here.
  \item Put derivations, alternative specifications, extra figures, and data details here.
  \item During Q\&A, jump directly to the relevant backup slide.
\end{itemize}
\end{frame}

\begin{frame}{Appendix: Example Derivation}
Suppose
\[
  \max_{c_1,c_2}\; u(c_1)+\beta u(c_2)
\]
subject to
\[
  c_1 + \frac{c_2}{1+r}=w.
\]
The first-order condition is
\[
  u'(c_1)=\beta(1+r)u'(c_2).
\]
\end{frame}

\begin{frame}{References}
\footnotesize
\begin{thebibliography}{9}
\bibitem{author2024}
Author, A. (2024). ``Paper Title.'' \textit{Journal Name}, 100(2), 1--25.

\bibitem{author2022}
Author, B. and C. Author (2022). ``Another Paper.'' \textit{Working Paper}.
\end{thebibliography}
\end{frame}

% ============================================================
\end{document}
